% GENERATED FILE. Edit canonical lessons, then run npm run generate:books.
% canonical-source-hash: fnv1a64:46f0b377534b65e4
% canonical-lessons: FA-C198-bayad, FA-C198-ide, FA-C198-farq, FA-C198-nemune, FA-C198-movaffaqiyat

\chapter{Must, Idea, Difference, Example, Success}
\label{ch:fa-must-idea-difference-example-success}
\hlchaptermodality{198}

\begin{chapteropening}
\textbf{By the end of this chapter:} \emph{I can say must, an idea, a difference, an example and success.}

Complete the last lesson of chapter 198: I can say must, an idea, a difference, an example and success.
\end{chapteropening}

\section[\texorpdfstring{\fa{باید} (bâyad)}{باید (bâyad)}]{\fa{باید} (bâyad) — must, should}

\label{lesson:FA-C198-bayad}



\begin{quote}
Before the new one: say the Persian for although, then the Persian for or.
\end{quote}

\subsection*{You'll want to know: \fa{باید}}
\textbf{\fa{باید}} — \emph{bâyad} — \textquotedblleft{}must, should\textquotedblright{}.

\begin{cousinweb}
\emph{bâyad} is a native Persian word, inherited rather than borrowed; it is usually traced to Middle Persian \emph{abāyēd}, \textquotedblleft{}it is necessary\textquotedblright{}.
\end{cousinweb}

\begin{grammarlens}[title={What you've built}]
One more word for this chapter.
\end{grammarlens}

\subsection*{Your turn}
\begin{itemize}
\raggedright
  \item \emph{Say it:} \emph{bâyad}
  \item \emph{Say it:} \emph{bâyad}, once more
  \item \emph{Say it:} say \emph{bâyad}
  \item \emph{Recall:} say the Persian for although, then the Persian for or, then say \emph{bâyad} again
\end{itemize}

\begin{tcolorbox}[breakable,colback=teal!4,colframe=teal!35!black,title={Before you move on}]
What does \fa{باید} mean? (\textquotedblleft{}Must, should\textquotedblright{}.) Say it once more.
\end{tcolorbox}
\section[\texorpdfstring{\fa{ایده} (ide)}{ایده (ide)}]{\fa{ایده} (ide) — an idea}

\label{lesson:FA-C198-ide}



\begin{quote}
Before the new one: say the Persian for or, then the Persian for must.
\end{quote}

\subsection*{You'll want to know: \fa{ایده}}
\textbf{\fa{ایده}} — \emph{ide} — \textquotedblleft{}an idea\textquotedblright{}.

\begin{cousinweb}
\emph{ide} is borrowed from French \emph{idée}, \textquotedblleft{}idea\textquotedblright{} (EtymDB). The French word goes back to Greek \emph{idea}, \textquotedblleft{}form, pattern\textquotedblright{}.
\end{cousinweb}

\begin{grammarlens}[title={What you've built}]
One more word for this chapter.
\end{grammarlens}

\subsection*{Your turn}
\begin{itemize}
\raggedright
  \item \emph{Say it:} \emph{ide}
  \item \emph{Say it:} \emph{ide}, once more
  \item \emph{Say it:} say \emph{ide}
  \item \emph{Recall:} say the Persian for or, then the Persian for must, then say \emph{ide} again
\end{itemize}

\begin{tcolorbox}[breakable,colback=teal!4,colframe=teal!35!black,title={Before you move on}]
What does \fa{ایده} mean? (\textquotedblleft{}An idea\textquotedblright{}.) Say it once more.
\end{tcolorbox}
\section[\texorpdfstring{\fa{فرق} (farq)}{فرق (farq)}]{\fa{فرق} (farq) — a difference}

\label{lesson:FA-C198-farq}



\begin{quote}
Before the new one: say the Persian for must, then the Persian for an idea.
\end{quote}

\subsection*{You'll want to know: \fa{فرق}}
\textbf{\fa{فرق}} — \emph{farq} — \textquotedblleft{}a difference\textquotedblright{}.

\begin{cousinweb}
\emph{farq} is borrowed from Arabic \emph{farq}, \textquotedblleft{}separation, difference\textquotedblright{}. Its Arabic root is \emph{f-r-q} (Lane).
\end{cousinweb}

\begin{grammarlens}[title={What you've built}]
One more word for this chapter.
\end{grammarlens}

\subsection*{Your turn}
\begin{itemize}
\raggedright
  \item \emph{Say it:} \emph{farq}
  \item \emph{Say it:} \emph{farq}, once more
  \item \emph{Say it:} say \emph{farq}
  \item \emph{Recall:} say the Persian for must, then the Persian for an idea, then say \emph{farq} again
\end{itemize}

\begin{tcolorbox}[breakable,colback=teal!4,colframe=teal!35!black,title={Before you move on}]
What does \fa{فرق} mean? (\textquotedblleft{}A difference\textquotedblright{}.) Say it once more.
\end{tcolorbox}
\section[\texorpdfstring{\fa{نمونه} (nemune)}{نمونه (nemune)}]{\fa{نمونه} (nemune) — an example, a sample}

\label{lesson:FA-C198-nemune}



\begin{quote}
Before the new one: say the Persian for an idea, then the Persian for a difference.
\end{quote}

\subsection*{You'll want to know: \fa{نمونه}}
\textbf{\fa{نمونه}} — \emph{nemune} — \textquotedblleft{}an example, a sample\textquotedblright{}.

\begin{cousinweb}
\emph{nemune} is built on \emph{nemudan}, \textquotedblleft{}to show\textquotedblright{}: what is shown. It is native, usually traced to Middle Persian \emph{nimūnag}.
\end{cousinweb}

\begin{grammarlens}[title={What you've built}]
One more word for this chapter.
\end{grammarlens}

\subsection*{Your turn}
\begin{itemize}
\raggedright
  \item \emph{Say it:} \emph{nemune}
  \item \emph{Say it:} \emph{nemune}, once more
  \item \emph{Say it:} say \emph{nemune}
  \item \emph{Recall:} say the Persian for an idea, then the Persian for a difference, then say \emph{nemune} again
\end{itemize}

\begin{tcolorbox}[breakable,colback=teal!4,colframe=teal!35!black,title={Before you move on}]
What does \fa{نمونه} mean? (\textquotedblleft{}An example, a sample\textquotedblright{}.) Say it once more.
\end{tcolorbox}
\section[\texorpdfstring{\fa{موفقیت} (movaffaqiyat)}{موفقیت (movaffaqiyat)}]{\fa{موفقیت} (movaffaqiyat) — success}

\label{lesson:FA-C198-movaffaqiyat}



\begin{quote}
Before the new one: say the Persian for a difference, then the Persian for an example.
\end{quote}

\subsection*{You'll want to know: \fa{موفقیت}}
\textbf{\fa{موفقیت}} — \emph{movaffaqiyat} — \textquotedblleft{}success\textquotedblright{}.

\begin{cousinweb}
\emph{movaffaqiyat} is borrowed from Arabic \emph{muwaffaqiyya}, \textquotedblleft{}success\textquotedblright{}. Its Arabic root is \emph{w-f-q} (Lane).
\end{cousinweb}

\begin{grammarlens}[title={What you've built}]
One more word for this chapter.
\end{grammarlens}

\subsection*{Your turn}
\begin{itemize}
\raggedright
  \item \emph{Say it:} \emph{movaffaqiyat}
  \item \emph{Say it:} \emph{movaffaqiyat}, once more
  \item \emph{Say it:} say \emph{movaffaqiyat}
  \item \emph{Recall:} say the Persian for a difference, then the Persian for an example, then say \emph{movaffaqiyat} again
\end{itemize}

\begin{tcolorbox}[breakable,colback=teal!4,colframe=teal!35!black,title={Before you move on}]
What does \fa{موفقیت} mean? (\textquotedblleft{}Success\textquotedblright{}.) Say it once more.
\end{tcolorbox}
